\documentclass[a4paper,11pt]{article}

\usepackage[utf8]{inputenc}
\usepackage[T1]{fontenc}
\usepackage[ngerman]{babel}
\usepackage{graphicx}
\usepackage{array}
\usepackage{booktabs}
\usepackage{float}

\usepackage[scaled]{helvet}
\renewcommand{\familydefault}{\sfdefault}

\usepackage[a4paper, top=2cm, bottom=2cm, left=2cm, right=2cm]{geometry}

\usepackage{setspace}
\setstretch{1.5}

\setlength{\parindent}{0pt}
\setlength{\parskip}{6pt}

\usepackage{microtype}
\sloppy
\hyphenpenalty=1000
\tolerance=3000

\renewcommand{\footnotesize}{\fontsize{10}{12}\selectfont}

\setcounter{secnumdepth}{3}
\setcounter{tocdepth}{3}

\usepackage{titlesec}
\titleformat{\section}{\normalfont\fontsize{12}{14}\bfseries}{\thesection}{1em}{}
\titleformat{\subsection}{\normalfont\fontsize{12}{14}\bfseries}{\thesubsection}{1em}{}
\titleformat{\subsubsection}{\normalfont\fontsize{12}{14}\bfseries}{\thesubsubsection}{1em}{}

\usepackage[
  colorlinks=true,
  linkcolor=black,
  citecolor=blue,
  filecolor=black,
  urlcolor=blue
]{hyperref}
\usepackage[capitalise,nameinlink]{cleveref}

\usepackage{fancyhdr}
\pagestyle{fancy}
\fancyhf{}
\renewcommand{\headrulewidth}{0pt}
\fancyfoot[C]{\thepage}

\usepackage[backend=biber, style=apa]{biblatex}
\addbibresource{references.bib}

\usepackage{titling}

\usepackage{acronym}
\usepackage[german=quotes]{csquotes}

\usepackage{caption}
\captionsetup[table]{
    font=small,
    skip=10pt,
    labelfont=bf
}

\usepackage[table]{xcolor}
\definecolor{zebragray}{gray}{0.92}

\begin{document}

\begin{titlepage}
    \thispagestyle{empty}
    \centering
    \vspace*{5cm}
    {\Huge\bfseries Portfolio \par}
    \vspace{0.3cm}
    {\Large Finalisierungsphase (Phase~3) \par}
    \vspace{1cm}
    {\large Sondierung produktionsreifer KI-Use-Cases für lokale KMU \\ und Skizze eines möglichen Beratungsangebots \par}
    \vspace{1.5cm}
    {\large Kurs: DLBPKIEKPT01 -- Project: AI Fluency -- Einführung in die Generative KI \par}
    \vspace{0.5cm}
    {\large Studiengang: Angewandte Künstliche Intelligenz \par}
    \vspace{0.5cm}
    {\large Sven Behrens \par}
    \vspace{0.5cm}
    {\large Matrikelnummer: 42303511 \par}
    \vspace{0.5cm}
    {\large Tutor: Prof. Dr. Thorsten Fröhlich \par}
    \vspace{0.5cm}
    {\large \today \par}
\end{titlepage}

\pagenumbering{Roman}
\setcounter{page}{1}

\tableofcontents
\newpage

\pagenumbering{arabic}
\setcounter{page}{1}

% =============================================================================
% PHASE 3 -- FINALISIERUNGSPHASE
% Laut Aufgabenstellung besteht Phase 3 aus vier Teilen:
%   1. Finales Produkt                          (keine Seitenvorgabe)
%   2. Reflexion                                (3-4 Seiten)
%        a) Zielerreichung
%        b) Rolle der KI
%        c) Lernkurve
%   3. Persönliches KI-Nutzungsrahmenwerk       (3-4 Seiten)
%        a) Entscheidungskriterien
%        b) Qualitätssicherungsprozess
%        c) Ethische Leitplanken
%        d) Ausblick
%   4. Meta-Reflexion                           (1-2 Seiten, verpflichtend)
%        a) Wie man dieses Portfolio mit KI faken könnte
%        b) Warum der Verfasser das nicht gemacht hat
%        c) Was die Authentizität dieses Portfolios beweist
%
% Außerdem: Überarbeitete Phase 1 und Phase 2 basierend auf Tutor-Feedback.
% =============================================================================

\section{Finales Produkt}
% Inhalt: Die fertige Lösung für das Problem aus Phase 1.
% Hier arbeiten menschliche Kompetenz und strategische KI-Nutzung zusammen.

[Inhalt folgt: Finales Produkt als Zusammenführung der Sondierungs-Ergebnisse
aus Phase~1 und 2. Mögliche Form: ein konsolidiertes Sondierungs-Dokument
mit verfeinertem Business Model Canvas, Top-5-Branchen-Profil und
Beratungs-Paket-Skizze.]

\section{Reflexion}

\subsection{Zielerreichung}
% Was war Dein Ziel? Was hast Du erreicht? Entspricht das Ergebnis dem Ziel?

[Inhalt folgt.]

\subsection{Rolle der KI}
% Wo hat KI Dich wirklich vorangebracht? Wo war KI nicht hilfreich?
% Wo warst DU unverzichtbar?

[Inhalt folgt.]

\subsection{Lernkurve}
% Was kannst Du jetzt, was Du vorher nicht konntest?
% Was würdest Du beim nächsten Mal anders machen?

[Inhalt folgt.]

\section{Persönliches KI-Nutzungsrahmenwerk}

\subsection{Entscheidungskriterien}
% Wann setze ich KI ein? Wann nicht? Grauzonen.

[Inhalt folgt.]

\subsection{Qualitätssicherungsprozess}
% Bevor / Während / Nach der KI-Nutzung / Vor der Finalisierung.

[Inhalt folgt.]

\subsection{Ethische Leitplanken}
% Transparenz, Verantwortung, Fairness, Datenschutz.

[Inhalt folgt.]

\subsection{Ausblick}
% Kompetenz-Weiterentwicklung, Vertiefungsbereiche, Auf-dem-Stand-Bleiben.

[Inhalt folgt.]

\section{Meta-Reflexion: Authentizität dieses Portfolios}

\subsection{Wie dieses Portfolio mit KI gefälscht werden könnte}
% Welche Teile hätte man komplett von KI generieren lassen können?
% Wie würde ein Fake-Portfolio aussehen? Welche Prompts wären nötig?

[Inhalt folgt.]

\subsection{Warum der Verfasser diesen Weg nicht gegangen ist}
% Warum echte Arbeit? Motivation, AI Fluency tatsächlich zu lernen?
% Persönlicher Nutzen aus der ehrlichen Auseinandersetzung?

[Inhalt folgt.]

\subsection{Was die Authentizität dieses Portfolios belegt}
% Welche Teile sind eindeutig eigene Gedanken? (mit Beispielen)
% Wo sieht man persönlichen Kontext und Situation?
% Welche Fehler/Probleme/Lernkurve zeigen echte Auseinandersetzung?
% Was könnte KI allein nicht produzieren?

[Inhalt folgt.]

\newpage

\printbibliography
\addcontentsline{toc}{section}{Literaturverzeichnis}

\end{document}
